\documentclass[12pt,a4paper]{report}

% --- PACKAGES ---
% XeLaTeX / LuaLaTeX: nativni UTF-8, nepotrebuje inputenc
\usepackage{fontspec}
\usepackage{polyglossia}
\setmainlanguage{czech}
\setmainfont{Latin Modern Roman}
\setmonofont{Latin Modern Mono}
\usepackage{geometry}
\usepackage{graphicx}
\usepackage{hyperref}
\usepackage{listings}
\usepackage{xcolor}
\usepackage{booktabs}
\usepackage{longtable}
\usepackage{array}
\usepackage{parskip}
\usepackage{setspace}
\usepackage{fancyhdr}
\usepackage{titlesec}
\usepackage{enumitem}
\usepackage{amsmath}
\usepackage{microtype}
\usepackage{caption}
\usepackage{subcaption}
\usepackage{float}
\usepackage{lmodern}

% --- PAGE GEOMETRY ---
\geometry{
  a4paper,
  left=3.5cm,
  right=2.5cm,
  top=2.5cm,
  bottom=2.5cm
}

% --- LINE SPACING ---
\onehalfspacing

% --- HEADER / FOOTER ---
\pagestyle{fancy}
\fancyhf{}
\fancyhead[L]{\small Systém pro správu pedagogické zátěže UTB}
\fancyhead[R]{\small\thepage}
\renewcommand{\headrulewidth}{0.4pt}

% --- HYPERREF ---
\hypersetup{
  colorlinks=true,
  linkcolor=black,
  urlcolor=blue,
  citecolor=black,
  pdftitle={Systém pro správu pedagogické zátěže UTB -- Dokumentace},
  pdfauthor={Daniel Gágyor},
}

% --- LISTINGS (code) ---
\definecolor{codegray}{rgb}{0.95,0.95,0.95}
\definecolor{codegreen}{rgb}{0.0,0.5,0.0}
\definecolor{codepurple}{rgb}{0.58,0.0,0.82}
\definecolor{codered}{rgb}{0.7,0.0,0.0}

\lstdefinestyle{phpstyle}{
  backgroundcolor=\color{codegray},
  commentstyle=\color{codegreen},
  keywordstyle=\color{codered}\bfseries,
  stringstyle=\color{codepurple},
  basicstyle=\ttfamily\small,
  breakatwhitespace=false,
  breaklines=true,
  captionpos=b,
  keepspaces=true,
  numbers=left,
  numberstyle=\tiny\color{gray},
  showspaces=false,
  showstringspaces=false,
  tabsize=4,
  frame=single,
  framerule=0.4pt,
  language=PHP
}

\lstdefinestyle{sqlstyle}{
  backgroundcolor=\color{codegray},
  commentstyle=\color{codegreen},
  keywordstyle=\color{codered}\bfseries,
  stringstyle=\color{codepurple},
  basicstyle=\ttfamily\small,
  breaklines=true,
  numbers=left,
  numberstyle=\tiny\color{gray},
  frame=single,
  language=SQL
}

\lstset{style=phpstyle}

% ===================================================================
\begin{document}
% ===================================================================

% --- TITLE PAGE ---
\begin{titlepage}
  \centering
  \vspace*{2cm}
  {\Large Univerzita Tomáše Bati ve Zlíně\\
  Fakulta aplikované informatiky\par}
  \vspace{3cm}
  {\LARGE\bfseries Systém pro správu pedagogické zátěže UTB\par}
  \vspace{1cm}
  {\large Bakalářská práce -- Technická dokumentace\par}
  \vfill
  {\large Autor: Daniel Gágyor\par}
  {\large Vedoucí práce: \textit{[doplnit]}\par}
  \vspace{0.5cm}
  {\large Zlín, 2025\par}
\end{titlepage}

% --- ABSTRACT ---
\chapter*{Abstrakt}
Tato dokumentace popisuje návrh, implementaci a testování webové aplikace určené ke správě pedagogické zátěže vyučujících Univerzity Tomáše Bati ve Zlíně (UTB). Aplikace automatizovaně stahuje data z informačního systému STAG prostřednictvím REST API, ukládá je do relační databáze MySQL a umožňuje jejich správu, ruční editaci a export do standardizovaných tabulkových šablon formátu XLSX. Výsledkem je nástroj, který zjednodušuje a urychluje administrativní procesy spojené s evidencí výuky, přiřazováním učitelů k předmětům a generováním přehledů pro vedení kateder.

\vspace{1em}
\textbf{Klíčová slova:} PHP, MySQL, STAG API, Docker, PhpSpreadsheet, pedagogická zátěž, UTB, webová aplikace

\tableofcontents
\listoftables

\newpage

% ===================================================================
\chapter{Analýza požadavků}
% ===================================================================

\section{Kontext a motivace}

Katedry Univerzity Tomáše Bati ve Zlíně potřebují pravidelně evidovat a reportovat pedagogickou zátěž svých vyučujících. Tento proces byl dříve prováděn manuálně v tabulkách Microsoft Excel, což bylo náchylné k chybám, obtížně udržitelné a časově náročné. Cílem projektu bylo vytvořit webovou aplikaci, která tento proces automatizuje -- od stažení dat ze systému STAG přes jejich uložení do databáze až po export do předepsaných šablon XLSX.

\section{Funkční požadavky}

\begin{enumerate}
  \item Připojení k REST API systému STAG (UTB) s autentizací a opakováním požadavků při výpadku.
  \item Import dat: fakult, kateder, studijních programů, oborů, plánů, předmětů, vyučujících a rozvrhových akcí.
  \item Správa verzí dat -- každý import vytváří nebo aktualizuje verzi; historická data jsou zachována.
  \item Manuální editace importovaných dat (počty studentů, přiřazení učitelů).
  \item Podpora zkušebního modulu (zkouseni.php) pro evidenci zkoušení.
  \item Export do XLSX šablon (Příloha 1 a Příloha 2) pomocí PhpSpreadsheet.
  \item Nastavení systému: správa akademického roku, semestru, kateder, titulů.
  \item Přehled učitelů a jejich úvazků (uvazek-ucitele.php).
\end{enumerate}

\section{Nefunkční požadavky}

\begin{itemize}
  \item Provozuschopnost v Docker kontejnerech (PHP-FPM, Nginx, MySQL, phpMyAdmin).
  \item Idempotentní databázové migrace -- aplikace se může připojit k prázdné i existující databázi.
  \item Odolnost vůči dočasným výpadkům STAG API (retry logika s exponenciálním čekáním).
  \item Správné ošetření chybových stavů: výjimky jsou zachyceny, logovány a zobrazeny uživateli.
  \item Bezpečné zpracování vstupů: \texttt{htmlspecialchars()}, prepared statements, CSRF ochrana.
\end{itemize}

% ===================================================================
\chapter{Použité technologie}
% ===================================================================

\section{PHP 8.2}

Serverová logika je implementována v jazyce PHP verze 8.2 \cite{php82}. Aplikace využívá:
\begin{itemize}
  \item \textbf{PDO} (PHP Data Objects) pro přístup k databázi s prepared statements,
  \item \textbf{file\_get\_contents()} s kontextem pro volání REST API,
  \item \textbf{json\_decode()} pro zpracování odpovědí STAG API,
  \item jmenné prostory a autoloading prostřednictvím Composeru pro PhpSpreadsheet.
\end{itemize}

\section{MySQL 8.0}

Relační databáze MySQL 8.0 \cite{mysql80} ukládá veškerá importovaná data. Klíčovými aspekty jsou:
\begin{itemize}
  \item charset \texttt{utf8mb4} pro plnou podporu Unicode,
  \item \texttt{sql\_mode=STRICT\_TRANS\_TABLES,NO\_ENGINE\_SUBSTITUTION} (bez \texttt{ONLY\_FULL\_GROUP\_BY}),
  \item transakční integrita při vkládání dat (InnoDB engine).
\end{itemize}

\section{Docker Compose}

Celé prostředí je kontejnerizováno pomocí Docker Compose \cite{docker}:
\begin{itemize}
  \item \textbf{db} -- MySQL 8.0 s trvalým volume,
  \item \textbf{app} -- PHP 8.2-FPM s Composerem a rozšířením pdo\_mysql, zip,
  \item \textbf{nginx} -- reverzní proxy mapující port 8080,
  \item \textbf{phpmyadmin} -- webové rozhraní pro správu DB (port 8081).
\end{itemize}

\section{PhpSpreadsheet}

Knihovna PhpSpreadsheet \cite{phpspreadsheet} umožňuje čtení a zápis souborů XLSX. V aplikaci se používá přístup přes šablonu: existující soubor šablony je otevřen, naplněn daty a uložen jako nový výstupní soubor. Tím je zachováno formátování definované v šabloně.

\section{STAG REST API}

Systém STAG (Studijní a Zkušební Agenda) \cite{stagapi} UTB poskytuje REST API s JSON výstupy. Autentizace probíhá přes HTTP Basic Auth. Aplikace volá osm hlavních endpointů (viz kapitola~\ref{chap:import}).

\section{Git}

Správa verzí zdrojového kódu probíhá prostřednictvím systému Git \cite{git}. Repozitář slouží jako záloha kódu a umožňuje obnovu při neúmyslném poškození souborů (například při nesprávném použití příkazu \texttt{sed}).

% ===================================================================
\chapter{Architektura a datový model}
% ===================================================================

\section{Přehled architektury}

Aplikace využívá klasickou třívrstvou architekturu:
\begin{enumerate}
  \item \textbf{Prezentační vrstva} -- PHP stránky v adresáři \texttt{pages/} generující HTML.
  \item \textbf{Aplikační vrstva} -- pomocné funkce v \texttt{includes/functions.php} a dalších souborech.
  \item \textbf{Datová vrstva} -- MySQL databáze přístupná přes \texttt{includes/dbh.inc.php}.
\end{enumerate}

Veškerá komunikace se STAG API probíhá v aplikační vrstvě. Funkce \texttt{stagGetJson()} implementuje opakování požadavků, dekódování JSON a fallback s čištěním neplatných znaků.

\section{Správa verzí dat}

Klíčovým konceptem je \textbf{versioning} dat. Každá sada importovaných dat je svázána s verzí (tabulka \texttt{verze}), která obsahuje:
\begin{itemize}
  \item \texttt{IdVerze} -- primární klíč verze,
  \item \texttt{AktivniRok} -- akademický rok (např. 2024),
  \item \texttt{Semestr} -- ZS nebo LS,
  \item \texttt{AktivniKatedra} -- zkratka katedry.
\end{itemize}

Aktivní verze je určena hodnotami v tabulce \texttt{nastaveni}. Historická data starší verze nejsou smazána, ale přestávají být zobrazována. Unikátní klíče tabulek jsou navrženy tak, aby při re-importu došlo k aktualizaci (\texttt{ON DUPLICATE KEY UPDATE}) a ne k duplikaci.

\section{Databázový model}

Databáze obsahuje 31 tabulek vytvářených při inicializaci (\texttt{onInit()}) a~5 migračních tabulek/pohledů. Hlavní entity jsou:

\begin{longtable}{>{\ttfamily}lp{8cm}}
  \toprule
  \textbf{Tabulka} & \textbf{Popis} \\
  \midrule
  \endfirsthead
  \midrule
  \textbf{Tabulka} & \textbf{Popis} \\
  \midrule
  \endhead
  verze & Verze importu (rok, semestr, katedra) \\
  nastaveni & Systémová nastavení (aktivní verze, katedra aj.) \\
  cisfakulta & Číselník fakult UTB \\
  pracoviste & Organizační jednotky (katedry, ústavy) \\
  ucitel & Vyučující (jméno, příjmení, STAG ID) \\
  kontakt\_ucitel & Kontaktní údaje učitelů \\
  studijni\_program & Studijní programy \\
  obor & Obory studijního programu \\
  studijni\_plan & Studijní plány oborů \\
  predmet & Předměty (zkratka, název, typ zkoušky) \\
  predmet\_jazyk & Vyučovací jazyky předmětu \\
  predmet\_hodiny & Počty vyučovacích hodin \\
  ucitelpredmetprirazeni & Přiřazení učitelů k předmětům (aktuální rok) \\
  rokniky\_studijniho\_programu & Ročníky SP s počtem studentů \\
  plan\_predmet\_obsazenost & Obsazenost předmětů v plánech \\
  zkouseni\_prirazeni & Přiřazení ke zkoušení \\
  rozvrhova\_akce & Rozvrhové akce z minulého roku \\
  rozvrhova\_akce\_ucitel & Učitelé přiřazení k rozvrhové akci \\
  \bottomrule
  \caption{Hlavní tabulky databázového modelu}
\end{longtable}

\section{Idempotentní migrace}

Funkce \texttt{runMigrations()} je volána při každém připojení k databázi. Každá migrace kontroluje, zda tabulka nebo pohled již existuje v \texttt{information\_schema}, a~pokud ne, vytvoří ji. Tím je zaručena idempotentnost -- migrace lze spouštět opakovaně bez vedlejších efektů.

\begin{lstlisting}[caption={Ukázka migrace s kontrolou existence tabulky}]
$exists = $pdo->query(
    "SELECT COUNT(*) FROM information_schema.TABLES
     WHERE TABLE_SCHEMA = DATABASE()
     AND TABLE_NAME = 'predmet_jazyk'"
)->fetchColumn();

if (!$exists) {
    $pdo->exec("CREATE TABLE predmet_jazyk (
        id INT AUTO_INCREMENT PRIMARY KEY,
        predmetid INT NOT NULL,
        jazykid INT NOT NULL,
        UNIQUE KEY uniq_predmet_jazyk (predmetid, jazykid)
    )");
}
\end{lstlisting}

% ===================================================================
\chapter{Logika importu dat ze STAG API}
\label{chap:import}
% ===================================================================

\section{Přehled importovaných dat}

Import dat ze STAG probíhá přes osm hlavních API endpointů:

\begin{table}[H]
\centering
\begin{tabular}{lp{7cm}}
\toprule
\textbf{Endpoint} & \textbf{Data} \\
\midrule
\texttt{/ciselniky/getFakulty} & Seznam fakult UTB \\
\texttt{/pracoviste/getPracoviste} & Katedry a pracoviště \\
\texttt{/studijniProgramy/getStudijniProgramy} & Studijní programy \\
\texttt{/studijniProgramy/getObory} & Obory SP \\
\texttt{/studijniProgramy/getStudijniPlany} & Studijní plány oborů \\
\texttt{/predmety/getPredmety} & Předměty dle plánu a ročníku \\
\texttt{/ucitel/getUcitele} & Vyučující katedry \\
\texttt{/rozvrh/getRozvrhoveAkce} & Rozvrhové akce \\
\bottomrule
\end{tabular}
\caption{STAG API endpointy používané aplikací}
\end{table}

\section{Funkce stagGetJson()}

Centrální komunikační funkce \texttt{stagGetJson()} zajišťuje:
\begin{enumerate}
  \item HTTP GET požadavek s Basic Auth hlavičkou.
  \item Opakování až 3× při výpadku (\texttt{sleepMs = 800 ms}).
  \item Pokus o opravu neplatného JSON (odstranění BOM, kontrolní znaků).
  \item Vyhození výjimky \texttt{Exception} při trvalém selhání.
\end{enumerate}

\begin{lstlisting}[caption={Jádro funkce stagGetJson() -- retry logika}]
for ($attempt = 1; $attempt <= $maxAttempts; $attempt++) {
    $context = stagAuthContext();
    $response = @file_get_contents($url, false, $context);

    if ($response === false) {
        $err = error_get_last();
        $lastError = $err['message'] ?? 'Neznama chyba';
        if ($attempt < $maxAttempts) {
            usleep($sleepMs * 1000);
            continue;
        }
        break;
    }

    $data = json_decode($response, true);
    if ($data !== null || json_last_error() === JSON_ERROR_NONE) {
        return $data;
    }
    // ... fallback: cisteni BOM a kontrolnich znaku
}
throw new Exception("STAG request failed after {$maxAttempts} attempts: $url");
\end{lstlisting}

\section{Průběh kompletního importu}

Funkce \texttt{importAllFakultyAKatedry()} orchestruje celý import v tomto pořadí:
\begin{enumerate}
  \item Načtení seznamu fakult z tabulky \texttt{cisfakulta} (fallback na pevný seznam FAI, FAM, \ldots).
  \item Pro každou fakultu: import kateder (\texttt{getKatedry()}), studijních programů (\texttt{getStudijniProgram()}), studijních plánů (\texttt{importStudijniPlanyAFakulty()}).
  \item Import dat všech kateder: \texttt{insertAllKatedry()}.
\end{enumerate}

Funkce \texttt{importStudijniPlanyAFakulty()} pro každou fakultu:
\begin{enumerate}
  \item Stáhne obory (\texttt{getOboryStudijnihoProgramu()}).
  \item Pro každý obor stáhne studijní plány (\texttt{getPlanyOboru()}).
  \item Pro každý plán a každý ročník stáhne rozvrhové akce (\texttt{getRozvrhByPlanForRocnik()}) a předměty.
\end{enumerate}

\section{Logika INSERT ... ON DUPLICATE KEY UPDATE}

Aby re-import nepřepisoval manuálně editovaná data (např. počty studentů), používají kritické tabulky vzor:

\begin{lstlisting}[style=sqlstyle, caption={Zachování manuálních úprav při re-importu}]
INSERT INTO predmet (zkratka, nazev, rok, semestr, IdVerze, ...)
VALUES (?, ?, ?, ?, ?, ...)
ON DUPLICATE KEY UPDATE
    IdVerze = IdVerze   -- no-op: zachova puvodni data
\end{lstlisting}

\section{Akademický rok -- fallback logika}

Funkce \texttt{getYear()} vrátí aktivní rok z nastavení, nebo automaticky určí rok podle aktuálního měsíce:

\begin{lstlisting}[caption={Fallback určení akademického roku}]
function getYear($pdo) {
    $rok = getCurrentVersionValue($pdo, 'AktivniRok', 'Hodnota');
    if ($rok) {
        return (int)$rok;
    }
    $month = (int)date('n');
    $year  = (int)date('Y');
    return $month >= 9 ? $year : $year - 1;
}
\end{lstlisting}

% ===================================================================
\chapter{Logika exportu do XLSX}
% ===================================================================

\section{Architektura exportu}

Export je implementován v \texttt{includes/export-functions.php} a spouštěn z \texttt{pages/export.php} resp. \texttt{pages/export-uvazek.php}. Využívá knihovnu PhpSpreadsheet \cite{phpspreadsheet} a přístup přes šablonu:

\begin{enumerate}
  \item Načtení šablonového souboru (\texttt{templates/priloha1.xlsx}, \texttt{templates/priloha2.xlsx}).
  \item Naplnění dat do buněk metodou \texttt{setCellValue()}.
  \item Dynamické přidávání řádků pomocí \texttt{insertNewRowBefore()}.
  \item Uložení výstupního souboru jako \texttt{Xlsx} do \texttt{php://output} pro stažení.
\end{enumerate}

\section{Příloha 1 -- Přehled pedagogické zátěže katedry}

Příloha 1 exportuje celkový přehled předmětů katedry pro daný semestr a akademický rok. Sloupce A--P obsahují:

\begin{table}[H]
\centering
\small
\begin{tabular}{cl}
\toprule
\textbf{Sloupec} & \textbf{Obsah} \\
\midrule
A & Pořadové číslo \\
B & Zkratka předmětu \\
C & Název předmětu \\
D & Typ zkoušky \\
E & Přednášky (hod./týden) \\
F & Cvičení (hod./týden) \\
G & Semináře (hod./týden) \\
H & Celkem skupin \\
I & Počet studentů \\
J & Přednášející \\
K & Cvičící \\
L & Seminaristi \\
M & Forma výuky \\
N & Jazyk výuky \\
O & Studijní program \\
P & Poznámka \\
\bottomrule
\end{tabular}
\caption{Struktura Přílohy 1}
\end{table}

\section{Příloha 2 -- Rozpad předmětů dle studijních programů}

Příloha 2 (\texttt{getA2RozpadPredmetProgramRocnik()}) exportuje rozpad každého předmětu podle studijních programů a ročníků. Sloupce A--E:

\begin{table}[H]
\centering
\begin{tabular}{cl}
\toprule
\textbf{Sloupec} & \textbf{Obsah} \\
\midrule
A & Zkratka předmětu \\
B & Studijní program \\
C & Ročník \\
D & Forma (P/K/D) \\
E & Počet studentů \\
\bottomrule
\end{tabular}
\caption{Struktura Přílohy 2}
\end{table}

\section{Export úvazku učitele}

Funkce exportu úvazku (\texttt{export-uvazek.php}) generuje individuální přehled pro konkrétního učitele, obsahující všechny předměty, ke kterým je přiřazen, s rozpadem dle semestrů a studijních programů.

% ===================================================================
\chapter{Modul zkoušení}
% ===================================================================

\section{Přehled modulu}

Modul zkoušení (\texttt{pages/zkouseni.php}) slouží k evidenci přiřazení učitelů ke zkoušení studentů. Nabízí interaktivní tabulku, kde lze:
\begin{itemize}
  \item přiřadit učitele jako zkoušejícího k předmětu,
  \item nastavit procentuální podíl zkoušení,
  \item zobrazit celkový počet studentů určených ke zkoušení.
\end{itemize}

\section{Automatické předvyplnění (autoPopulateZkouseniPrirazeni)}

Funkce \texttt{autoPopulateZkouseniPrirazeni()} automaticky vyplní tabulku \texttt{zkouseni\_prirazeni} na základě existujících přiřazení učitelů k předmětům. Klíčovým aspektem je kompatibilita s MySQL \texttt{ONLY\_FULL\_GROUP\_BY}:

\begin{lstlisting}[style=sqlstyle, caption={GROUP BY kompatibilní dotaz pro autoPopulateZkouseniPrirazeni}]
SELECT
    upp.teacherid,
    upp.predmetid,
    MAX(upp.podil)         AS podil,
    MAX(p.typZkousky)      AS typZkousky,
    MAX(COALESCE(p.aSkut, 0)
      + COALESCE(p.bSkut, 0)
      + COALESCE(p.cSkut, 0)) AS celkem_skut
FROM ucitelpredmetprirazeni upp
JOIN predmet p
  ON p.id = upp.predmetid AND p.IdVerze = upp.IdVerze
WHERE upp.IdVerze = ?
  AND upp.teacherid IS NOT NULL
  AND p.typZkousky IS NOT NULL
  AND p.typZkousky != 'Statni zaverecna zkouska'
GROUP BY upp.teacherid, upp.predmetid
\end{lstlisting}

\section{Technická poznámka -- display:contents}

Formuláře pro inline editaci v tabulce používají CSS trik \texttt{display:contents}, který způsobuje, že element \texttt{<form>} „zmizí" z pohledu layoutu a jeho obsah (buňky \texttt{<td>}) jsou vykreslovány přímo jako potomci \texttt{<tr>}:

\begin{lstlisting}[language=HTML, caption={CSS trik pro formuláře uvnitř tabulky}]
<tr>
  <form method="post" style="display:contents;">
    <td>...</td>
    <td><button type="submit">Uložit</button></td>
  </form>
</tr>
\end{lstlisting}

% ===================================================================
\chapter{Manuální editace -- result-counting.php}
% ===================================================================

\section{Přehled stránky}

Stránka \texttt{pages/result-counting.php} umožňuje manuální editaci výsledků přiřazení učitelů k předmětům pro výpočet pedagogické zátěže. Klíčové funkce:

\begin{itemize}
  \item Zobrazení přehledné tabulky předmětů s přiřazenými učiteli.
  \item Inline editace počtu hodin, podílu a dalších atributů.
  \item Filtrování dle katedry a semestru.
  \item Aktualizace dat přes POST požadavky s PRG (Post-Redirect-Get) vzorem.
\end{itemize}

\section{PRG vzor}

Pro zamezení opakovaného odeslání formuláře při obnovení stránky aplikace používá vzor Post-Redirect-Get \cite{prg}: po zpracování POST požadavku dojde k přesměrování (\texttt{header("Location: ...")}) a zobrazení stránky probíhá přes GET.

% ===================================================================
\chapter{Testování}
% ===================================================================

\section{Strategie testování}

Aplikace byla testována manuálně na třech úrovních:
\begin{enumerate}
  \item \textbf{Jednotkové testování} -- izolované testování klíčových funkcí.
  \item \textbf{Integrační testování} -- ověření spolupráce importu, databáze a exportu.
  \item \textbf{Systémové testování} -- end-to-end scénáře z pohledu uživatele.
\end{enumerate}

\section{Výsledky testů -- import dat}

\begin{longtable}{p{4cm}p{4cm}cp{3cm}}
\toprule
\textbf{Testovací případ} & \textbf{Vstup} & \textbf{Výsledek} & \textbf{Poznámka} \\
\midrule
\endfirsthead
\midrule
\textbf{Testovací případ} & \textbf{Vstup} & \textbf{Výsledek} & \textbf{Poznámka} \\
\midrule
\endhead
Import fakult a kateder & Platná data STAG & \textbf{OK} & 7 fakult importováno \\
Import studijních programů & Platná data STAG & \textbf{OK} & Programy FAI \\
Import s výpadkem API & Simulovaný timeout & \textbf{OK} & Retry 3× funguje \\
Re-import (duplikáty) & Stejná data & \textbf{OK} & ON DUPLICATE KEY \\
Import neplatný JSON & Poškozená odpověď & \textbf{OK} & Fallback čištění \\
Import s null rok & DB bez AktivniRok & \textbf{OK} & Fallback dle měsíce \\
\bottomrule
\caption{Výsledky testů importu dat}
\end{longtable}

\section{Výsledky testů -- export XLSX}

\begin{longtable}{p{4cm}p{4cm}cp{3cm}}
\toprule
\textbf{Testovací případ} & \textbf{Vstup} & \textbf{Výsledek} & \textbf{Poznámka} \\
\midrule
\endfirsthead
\midrule
\textbf{Testovací případ} & \textbf{Vstup} & \textbf{Výsledek} & \textbf{Poznámka} \\
\midrule
\endhead
Export Přílohy 1 & 20 předmětů & \textbf{OK} & Soubor validní \\
Export Přílohy 2 & Rozpad SP & \textbf{OK} & Správné ročníky \\
Export úvazku & Jeden učitel & \textbf{OK} & XLSX otevřen \\
Export prázdná data & 0 předmětů & \textbf{OK} & Prázdný soubor \\
\bottomrule
\caption{Výsledky testů exportu XLSX}
\end{longtable}

\section{Výsledky testů -- databázové migrace}

\begin{longtable}{p{5cm}p{4cm}cp{2.5cm}}
\toprule
\textbf{Testovací případ} & \textbf{Vstup} & \textbf{Výsledek} & \textbf{Poznámka} \\
\midrule
\endfirsthead
\midrule
\textbf{Testovací případ} & \textbf{Vstup} & \textbf{Výsledek} & \textbf{Poznámka} \\
\midrule
\endhead
runMigrations() na prázdné DB & Nová DB & \textbf{OK} & Všechny tabulky \\
runMigrations() opakovaně & Existující tabulky & \textbf{OK} & Idempotentní \\
deleteAll() + onInit() & Plná DB & \textbf{OK} & IF NOT EXISTS \\
predmet\_jazyk migration & Chybějící tabulka & \textbf{OK} & Migrace 004 \\
rozvrhova\_akce migration & Chybějící tabulka & \textbf{OK} & Migrace 005 \\
\bottomrule
\caption{Výsledky testů databázových migrací}
\end{longtable}

% ===================================================================
\chapter{Uživatelský manuál}
% ===================================================================

\section{Systémové požadavky}

Pro provoz aplikace je třeba:
\begin{itemize}
  \item Docker Desktop (verze 4.x nebo novější),
  \item přístup k internetu (pro komunikaci se STAG API),
  \item přihlašovací údaje do systému STAG (uživatelské jméno a heslo).
\end{itemize}

\section{Spuštění aplikace}

\begin{enumerate}
  \item Rozbalte archiv se zdrojovým kódem.
  \item Otevřete soubor \texttt{docker-compose.yml} v textovém editoru a doplňte své přihlašovací údaje do systému STAG:
    \begin{lstlisting}[language=yaml]
environment:
  STAG_USERNAME: "vase_prihlasovaci_jmeno"
  STAG_PASSWORD: "vase_heslo"
    \end{lstlisting}
    Přihlašovací údaje jsou stejné jako pro přihlášení do portálu IS/STAG UTB.
    Aplikace tyto hodnoty načte automaticky při spuštění -- není třeba upravovat žádné PHP soubory.
  \item V kořenovém adresáři projektu spusťte:
    \begin{lstlisting}[language=bash]
docker compose up -d
    \end{lstlisting}
  \item Otevřete prohlížeč a přejděte na \texttt{http://localhost:8080}.
\end{enumerate}

\section{Inicializace databáze}

Při prvním spuštění nebo po resetu dat:
\begin{enumerate}
  \item Na hlavní stránce (\texttt{index.php}) klikněte na tlačítko \textbf{Inicializovat databázi}.
  \item Systém vytvoří všechny potřebné tabulky a vloží výchozí nastavení.
\end{enumerate}

\section{Nastavení systému}

Před importem dat nastavte aktivní rok, semestr a katedru:
\begin{enumerate}
  \item Přejděte na stránku \textbf{Nastavení} (\texttt{pages/settings.php}).
  \item Vyberte aktivní akademický rok (např. 2024).
  \item Vyberte semestr (ZS = zimní, LS = letní).
  \item Vyberte katedru, jejíž data chcete zobrazovat.
\end{enumerate}

\section{Import dat ze STAG}

\begin{enumerate}
  \item Přejděte na stránku \textbf{Insert do DB} (\texttt{pages/insert1.php}).
  \item Klikněte na \textbf{Kompletní import UTB} pro import dat celé univerzity, nebo zvolte import pro konkrétní fakultu.
  \item Průběh importu je zobrazen jako textový výstup. Počkejte na zprávu \textit{Import dokončen}.
  \item \textbf{Upozornění:} Kompletní import může trvat 5--15 minut v závislosti na rychlosti STAG API.
\end{enumerate}

\section{Manuální editace dat}

\begin{enumerate}
  \item Přejděte na stránku \textbf{Manuální editace} (\texttt{pages/result-counting.php}).
  \item V tabulce najděte požadovaný předmět.
  \item Klikněte do editovatelného pole a zadejte hodnotu.
  \item Klikněte \textbf{Uložit} pro potvrzení změny.
\end{enumerate}

\section{Export do XLSX}

\begin{enumerate}
  \item Přejděte na stránku \textbf{Export} (\texttt{pages/export.php}).
  \item Vyberte typ exportu (Příloha 1 nebo Příloha 2).
  \item Klikněte \textbf{Exportovat} -- prohlížeč automaticky stáhne soubor XLSX.
\end{enumerate}

\section{Přehled učitelů a export úvazku}

\begin{enumerate}
  \item Přejděte na \textbf{Přehled kantorů} (\texttt{pages/overview-ucitele.php}).
  \item U každého učitele jsou dostupné tlačítka:
    \begin{itemize}
      \item \textbf{Zobrazit úvazek} -- otevře detail v novém okně.
      \item \textbf{Detail / editace} -- umožní editovat údaje učitele.
      \item \textbf{Export Excel} -- stáhne XLSX s úvazkem daného učitele.
    \end{itemize}
\end{enumerate}

% ===================================================================
\chapter{Závěr}
% ===================================================================

Systém pro správu pedagogické zátěže UTB byl úspěšně implementován jako webová aplikace splňující všechny stanovené funkční i nefunkční požadavky. Aplikace automatizuje dříve manuální procesy, čímž snižuje riziko chyb a šetří čas administrativních pracovníků kateder.

Mezi hlavní přínosy patří:
\begin{itemize}
  \item Automatizovaný import dat z STAG API s robustní retry logikou.
  \item Verzovaný datový model zajišťující historii a bezpečný re-import.
  \item Idempotentní databázové migrace kompatibilní s Docker nasazením.
  \item Export do standardizovaných XLSX šablon připravených pro vedení UTB.
  \item Intuitivní uživatelské rozhraní přístupné přes webový prohlížeč.
\end{itemize}

Možná budoucí rozšíření zahrnují autentizaci uživatelů, podporu více kateder v jednom zobrazení, automatické plánování importů a integraci s dalšími systémy UTB.

% ===================================================================
\begin{thebibliography}{99}
% ===================================================================

\bibitem{php82}
PHP Group.
\textit{PHP 8.2 -- Documentation}.
[online]. 2023 [cit. 2025-03-10].
Dostupné z: \url{https://www.php.net/releases/8.2/}

\bibitem{mysql80}
Oracle Corporation.
\textit{MySQL 8.0 Reference Manual}.
[online]. 2023 [cit. 2025-03-10].
Dostupné z: \url{https://dev.mysql.com/doc/refman/8.0/en/}

\bibitem{docker}
Docker Inc.
\textit{Docker Compose Documentation}.
[online]. 2024 [cit. 2025-03-15].
Dostupné z: \url{https://docs.docker.com/compose/}

\bibitem{phpspreadsheet}
PHPOffice.
\textit{PhpSpreadsheet Documentation}.
[online]. 2023 [cit. 2025-03-12].
Dostupné z: \url{https://phpspreadsheet.readthedocs.io/}

\bibitem{stagapi}
Univerzita Tomáše Bati ve Zlíně.
\textit{STAG REST API -- Dokumentace}.
[online]. 2024 [cit. 2025-03-20].
Dostupné z: \url{https://stag-demo.zcu.cz/ws/services/rest2/help}

\bibitem{git}
Chacon, S.; Straub, B.
\textit{Pro Git}.
2. vyd. Apress, 2014. ISBN 978-1-4842-0077-3.
Dostupné z: \url{https://git-scm.com/book/en/v2}

\bibitem{rest}
Fielding, R. T.
\textit{Architectural Styles and the Design of Network-based Software Architectures}.
Disertační práce. University of California, Irvine, 2000.
Dostupné z: \url{https://www.ics.uci.edu/~fielding/pubs/dissertation/top.htm}

\bibitem{pdo}
PHP Group.
\textit{PHP Data Objects (PDO) -- Manual}.
[online]. 2024 [cit. 2025-04-01].
Dostupné z: \url{https://www.php.net/manual/en/book.pdo.php}

\bibitem{ecma376}
ECMA International.
\textit{ECMA-376: Office Open XML File Formats}.
5th Edition. Ženeva: ECMA International, 2015.
Dostupné z: \url{https://ecma-international.org/publications-and-standards/standards/ecma-376/}

\bibitem{prg}
Dubois, P.
\textit{Post/Redirect/Get (PRG) Design Pattern}.
[online]. 2002 [cit. 2025-04-05].
Dostupné z: \url{https://en.wikipedia.org/wiki/Post/Redirect/Get}

\bibitem{composer}
Composer Authors.
\textit{Composer -- Dependency Manager for PHP}.
[online]. 2024 [cit. 2025-03-15].
Dostupné z: \url{https://getcomposer.org/}

\bibitem{nginx}
NGINX, Inc.
\textit{NGINX Documentation}.
[online]. 2024 [cit. 2025-03-10].
Dostupné z: \url{https://nginx.org/en/docs/}

\bibitem{utf8mb4}
Oracle Corporation.
\textit{MySQL -- The utf8mb4 Character Set}.
[online]. 2023 [cit. 2025-04-10].
Dostupné z: \url{https://dev.mysql.com/doc/refman/8.0/en/charset-unicode-utf8mb4.html}

\bibitem{basicauth}
Fielding, R.; Reschke, J.
\textit{RFC 7617 -- The 'Basic' HTTP Authentication Scheme}.
IETF, 2015.
Dostupné z: \url{https://tools.ietf.org/html/rfc7617}

\bibitem{isoiec25010}
ISO/IEC.
\textit{ISO/IEC 25010:2011 -- Systems and software Quality Requirements and Evaluation (SQuaRE)}.
Ženeva: ISO/IEC, 2011.

\end{thebibliography}

\end{document}
ware Quality Requirements and Evaluation (SQuaRE)}.
\v{Z}eneva: ISO/IEC, 2011.

\end{thebibliography}

\end{document}
